\scanpage{76}
d) Il existe une décomposition de $\deleted{\operatorname{Hom}}P(X,Y)$
en
\[
P(X,Y)=\bigoplus_{n\geq0}P_n(X,Y)
\]
telle que pour $f_n\in P_n(X,Y)$, et $\lambda\in\mathbb Z$, on ait
\[
f_n\circ(\lambda\mathrm{id}_X)=\lambda^nf_n \tag*{$*$}
\]
Plus généralement, si $X_1,\ldots,X_r,Y$ sont des
objets de $P$, on a une décomposition
\[
P(\deleted{\unsure{\prod_i X_i}}X_1\times\cdots\times X_r,Y)=
\bigoplus_{n=(n_1,\ldots,n_r)\in\mathbb N^r}P_n(X_1,\ldots,X_r;Y)\tag*{($*$)}
\]
avec, pour $f\in P_{n_1,\ldots,n_r}(X_1,\ldots,X_r;Y)$ et $\lambda=(\lambda_i)_{1\leq i\leq r}\in\mathbb Z^r$,
on ait
\[
f\circ(\lambda_1\mathrm{id}_{X_1},\ldots,\lambda_r\mathrm{id}_{X_r})
=\lambda_1^{n_1}\lambda_r^{n_r}f.
\]

% margin note written sideways
\margin{\uline{Ex} si $P$ est un groupe annulé par $p$, alors l'op. de $\mathbb Z$ par $\lambda\longmapsto\lambda\,\mathrm{id}_P$ et $\lambda\longmapsto\lambda^p\,\mathrm{id}_{\mathbb Z}$ est la $\underline{=}$, donc on peut concentrer au choix $P$ en \unsure{un|une} \unsure{en} degré $1$, ou $p$ !}

\uline{Remarques} Cette décomposition \inserted{n'}est \inserted{pas} \unsure{unique}, un
\deleted{\unsure{sorte de|retour de|vertu de}} canular ``arithmétique'' -- il le serait

\deleted{\uline{Lemme} Soit $P$ un \uline{groupe abélien}, et
\[
u:(\mathbb Z^\times)^r\longrightarrow\operatorname{End}(P)
\]
un hom de monoïdes.
\deleted{Pour tt $n=(n_1,\ldots,n_r)\in\mathbb Z$,}
Supposons qu'il existe une décomposition en
somme directe
\[
P=\sum_{n=(n_1,\ldots,n_r)\in\mathbb N^r} P_n\unsure{,}
\]
telle que pour $f\in P_n$\inserted{ et $\lambda=(\lambda_1,\ldots,\lambda_r)\in\mathbb Z^r$,} on ait
\[
\deleted{\unsure{\mu_1}}u(\lambda)\cdot f=\lambda_1^{n_1}\cdots\lambda_r^{n_r}f.
\]
Alors cette décomposition est unique, plus
précisément on a pour tt $n\in\mathbb N^r$
\[
P_n=\{f\in P\mid u(\lambda)f=\lambda_1^{n_1}\cdots\lambda_r^{n_r}f\quad\forall\lambda\in\mathbb Z^r\}
\]
Il suffit de voir que si $n'\ne n$ et si
$g_{n'}\in P_{n'}$ est tel que $\unsure{\mu|u}(\lambda)g_{n'}=\lambda_1^{n_1}\cdots\lambda_r^{n_r}g_{n'}\quad\forall\lambda\in\mathbb Z^r$,
alors $g_{n'}=0$. Or on aura aussi
\[
\deleted{\unsure{\lambda}}u(\lambda)g_{n'}=\lambda_1^{n'_1}\cdots\lambda_r^{n'_r}g_{n'}
\]}
% margin note written sideways; outer right edge cut by the scan
\margin{\unsure{Il} \unsure{si} \illegible{1}}
% archive sheet number: 75
